\subsection{模型价值的证明性讨论}
模型的有用性不仅取决于是否能够解释一个例子，还取决于它能否在新环境中给出稳定预测。准确性要求预测与观测相符；广泛性要求同一机制可迁移；简洁性要求不以无约束参数堆积换取拟合；识别性要求参数可由数据区分；可证伪性要求存在可能否定模型的事实。六者可能发生张力，研究者应在文中明确取舍。

\subsection{从反例学习}
反例不是理论失败的附录，而是界定适用范围的工具。若一个命题在非闭预算集上失败，应修正存在性条件；若在非凸偏好上失败，应修正唯一性结论；若在不可比关系上失败，应放弃单一序数效用的强断言。每个反例都对应一条必须补充的假设。

\subsection{综合案例：平台规则改变}
设平台规则使可行集从 $B_0$ 变为 $B_1$。正向问题是比较 $C(B_0)$ 与 $C(B_1)$；规范问题则要进一步比较效用、收入分配与外部性。不能把选择变化直接等同于福利变化，因为可行集扩张可能同时伴随风险上升或信息成本增加。

\subsection{本节练习}
\begin{enumerate}
\item 给出一个可行集扩张但最大值下降的例子，并说明需要什么效用定义。
\item 说明为什么选择对应比选择点更适合表示多重最优。
\item 对一个政策反事实分别写出正向命题、机制命题和规范命题。
\end{enumerate}
